% Options for packages loaded elsewhere
\PassOptionsToPackage{unicode}{hyperref}
\PassOptionsToPackage{hyphens}{url}
\documentclass[
]{article}
\usepackage{xcolor}
\usepackage[margin=1in]{geometry}
\usepackage{amsmath,amssymb}
\setcounter{secnumdepth}{-\maxdimen} % remove section numbering
\usepackage{iftex}
\ifPDFTeX
  \usepackage[T1]{fontenc}
  \usepackage[utf8]{inputenc}
  \usepackage{textcomp} % provide euro and other symbols
\else % if luatex or xetex
  \usepackage{unicode-math} % this also loads fontspec
  \defaultfontfeatures{Scale=MatchLowercase}
  \defaultfontfeatures[\rmfamily]{Ligatures=TeX,Scale=1}
\fi
\usepackage{lmodern}
\ifPDFTeX\else
  % xetex/luatex font selection
\fi
% Use upquote if available, for straight quotes in verbatim environments
\IfFileExists{upquote.sty}{\usepackage{upquote}}{}
\IfFileExists{microtype.sty}{% use microtype if available
  \usepackage[]{microtype}
  \UseMicrotypeSet[protrusion]{basicmath} % disable protrusion for tt fonts
}{}
\makeatletter
\@ifundefined{KOMAClassName}{% if non-KOMA class
  \IfFileExists{parskip.sty}{%
    \usepackage{parskip}
  }{% else
    \setlength{\parindent}{0pt}
    \setlength{\parskip}{6pt plus 2pt minus 1pt}}
}{% if KOMA class
  \KOMAoptions{parskip=half}}
\makeatother
\usepackage{color}
\usepackage{fancyvrb}
\newcommand{\VerbBar}{|}
\newcommand{\VERB}{\Verb[commandchars=\\\{\}]}
\DefineVerbatimEnvironment{Highlighting}{Verbatim}{commandchars=\\\{\}}
% Add ',fontsize=\small' for more characters per line
\usepackage{framed}
\definecolor{shadecolor}{RGB}{248,248,248}
\newenvironment{Shaded}{\begin{snugshade}}{\end{snugshade}}
\newcommand{\AlertTok}[1]{\textcolor[rgb]{0.94,0.16,0.16}{#1}}
\newcommand{\AnnotationTok}[1]{\textcolor[rgb]{0.56,0.35,0.01}{\textbf{\textit{#1}}}}
\newcommand{\AttributeTok}[1]{\textcolor[rgb]{0.13,0.29,0.53}{#1}}
\newcommand{\BaseNTok}[1]{\textcolor[rgb]{0.00,0.00,0.81}{#1}}
\newcommand{\BuiltInTok}[1]{#1}
\newcommand{\CharTok}[1]{\textcolor[rgb]{0.31,0.60,0.02}{#1}}
\newcommand{\CommentTok}[1]{\textcolor[rgb]{0.56,0.35,0.01}{\textit{#1}}}
\newcommand{\CommentVarTok}[1]{\textcolor[rgb]{0.56,0.35,0.01}{\textbf{\textit{#1}}}}
\newcommand{\ConstantTok}[1]{\textcolor[rgb]{0.56,0.35,0.01}{#1}}
\newcommand{\ControlFlowTok}[1]{\textcolor[rgb]{0.13,0.29,0.53}{\textbf{#1}}}
\newcommand{\DataTypeTok}[1]{\textcolor[rgb]{0.13,0.29,0.53}{#1}}
\newcommand{\DecValTok}[1]{\textcolor[rgb]{0.00,0.00,0.81}{#1}}
\newcommand{\DocumentationTok}[1]{\textcolor[rgb]{0.56,0.35,0.01}{\textbf{\textit{#1}}}}
\newcommand{\ErrorTok}[1]{\textcolor[rgb]{0.64,0.00,0.00}{\textbf{#1}}}
\newcommand{\ExtensionTok}[1]{#1}
\newcommand{\FloatTok}[1]{\textcolor[rgb]{0.00,0.00,0.81}{#1}}
\newcommand{\FunctionTok}[1]{\textcolor[rgb]{0.13,0.29,0.53}{\textbf{#1}}}
\newcommand{\ImportTok}[1]{#1}
\newcommand{\InformationTok}[1]{\textcolor[rgb]{0.56,0.35,0.01}{\textbf{\textit{#1}}}}
\newcommand{\KeywordTok}[1]{\textcolor[rgb]{0.13,0.29,0.53}{\textbf{#1}}}
\newcommand{\NormalTok}[1]{#1}
\newcommand{\OperatorTok}[1]{\textcolor[rgb]{0.81,0.36,0.00}{\textbf{#1}}}
\newcommand{\OtherTok}[1]{\textcolor[rgb]{0.56,0.35,0.01}{#1}}
\newcommand{\PreprocessorTok}[1]{\textcolor[rgb]{0.56,0.35,0.01}{\textit{#1}}}
\newcommand{\RegionMarkerTok}[1]{#1}
\newcommand{\SpecialCharTok}[1]{\textcolor[rgb]{0.81,0.36,0.00}{\textbf{#1}}}
\newcommand{\SpecialStringTok}[1]{\textcolor[rgb]{0.31,0.60,0.02}{#1}}
\newcommand{\StringTok}[1]{\textcolor[rgb]{0.31,0.60,0.02}{#1}}
\newcommand{\VariableTok}[1]{\textcolor[rgb]{0.00,0.00,0.00}{#1}}
\newcommand{\VerbatimStringTok}[1]{\textcolor[rgb]{0.31,0.60,0.02}{#1}}
\newcommand{\WarningTok}[1]{\textcolor[rgb]{0.56,0.35,0.01}{\textbf{\textit{#1}}}}
\usepackage{graphicx}
\makeatletter
\newsavebox\pandoc@box
\newcommand*\pandocbounded[1]{% scales image to fit in text height/width
  \sbox\pandoc@box{#1}%
  \Gscale@div\@tempa{\textheight}{\dimexpr\ht\pandoc@box+\dp\pandoc@box\relax}%
  \Gscale@div\@tempb{\linewidth}{\wd\pandoc@box}%
  \ifdim\@tempb\p@<\@tempa\p@\let\@tempa\@tempb\fi% select the smaller of both
  \ifdim\@tempa\p@<\p@\scalebox{\@tempa}{\usebox\pandoc@box}%
  \else\usebox{\pandoc@box}%
  \fi%
}
% Set default figure placement to htbp
\def\fps@figure{htbp}
\makeatother
\setlength{\emergencystretch}{3em} % prevent overfull lines
\providecommand{\tightlist}{%
  \setlength{\itemsep}{0pt}\setlength{\parskip}{0pt}}
\usepackage{bookmark}
\IfFileExists{xurl.sty}{\usepackage{xurl}}{} % add URL line breaks if available
\urlstyle{same}
\hypersetup{
  hidelinks,
  pdfcreator={LaTeX via pandoc}}

\author{}
\date{\vspace{-2.5em}}

\begin{document}

\section{Analysis Notes}\label{analysis-notes}

\begin{itemize}
\tightlist
\item
  BR = bark roughness
\item
  CT = bark condensed tannins
\item
  CN = bark carbon-nitrogen ratio
\item
  PC = percent cover (percent occupied cells in the sampling grid)
\item
  SR = lichen species richness
\item
  L = size of lichen network (number of nodes)
\item
  Cen = centrality of the lichen network
\end{itemize}

\begin{Shaded}
\begin{Highlighting}[]
\DocumentationTok{\#\# Make sure to run make.R first and load the necessary objects.}

\FunctionTok{library}\NormalTok{(pacman)}
\FunctionTok{p\_load}\NormalTok{(drake, xtable, rmarkdown, flextable, knitr)}
\FunctionTok{loadd}\NormalTok{(onc.dat)}
\FunctionTok{loadd}\NormalTok{(cn.d.onc)}
\end{Highlighting}
\end{Shaded}

\subsection{Is the network similarity driven by abundance of
lichen?}\label{is-the-network-similarity-driven-by-abundance-of-lichen}

We found support for the hypothesis that genotypically based variation
in lichen network structure is potentially driven by variation in bark
roughness, not purely as a function of increasing lichen abundance or
richness. In an initital PerMANOVA with sequential partitioning of
variance explained, genotype has a strong, statistically significant
effect when including tree traits and lichen community indices in the
model. In this ordering, condensed tannins and species richness are both
significant, although with low explanatory power. After moving genotype
to the final predictor in the model, genotype is no longer significant
and has low explanatory power, while both bark roughness and species
richness are now significant and have increased in the variance in
network similarity explained. Taken together, these results suggest that
there is potentially a causal relationship between genotypically
explained variance in bark roughness and lichen community species
richness that affects the similarity of lichen network structure.

\begin{Shaded}
\begin{Highlighting}[]
\NormalTok{cn.geno.trait.pc.sr }\OtherTok{\textless{}{-}}\NormalTok{ vegan}\SpecialCharTok{::}\FunctionTok{adonis2}\NormalTok{(}
\NormalTok{                               cn.d.onc}\SpecialCharTok{\textasciitilde{}}\NormalTok{ geno }\SpecialCharTok{+}\NormalTok{ CT }\SpecialCharTok{+}\NormalTok{ pH }\SpecialCharTok{+}\NormalTok{ CN }\SpecialCharTok{+}\NormalTok{ BR }\SpecialCharTok{+}\NormalTok{ PC }\SpecialCharTok{+}\NormalTok{ SR,}
                               \AttributeTok{by =} \StringTok{"term"}\NormalTok{, }
                               \AttributeTok{data =}\NormalTok{ onc.dat, }
                               \AttributeTok{mrank =} \ConstantTok{TRUE}\NormalTok{,}
                               \AttributeTok{permutations =} \DecValTok{100000}\NormalTok{)}

\NormalTok{cn.trait.pc.sr.geno  }\OtherTok{\textless{}{-}}\NormalTok{ vegan}\SpecialCharTok{::}\FunctionTok{adonis2}\NormalTok{(}
\NormalTok{                                   cn.d.onc}\SpecialCharTok{\textasciitilde{}}\NormalTok{ CT }\SpecialCharTok{+}\NormalTok{ pH }\SpecialCharTok{+}\NormalTok{ CN }\SpecialCharTok{+}\NormalTok{ BR }\SpecialCharTok{+}\NormalTok{ PC }\SpecialCharTok{+}\NormalTok{ SR }\SpecialCharTok{+}\NormalTok{ geno,}
                                   \AttributeTok{by =} \StringTok{"term"}\NormalTok{, }
                                   \AttributeTok{data =}\NormalTok{ onc.dat, }
                                   \AttributeTok{mrank =} \ConstantTok{TRUE}\NormalTok{,}
                                   \AttributeTok{permutations =} \DecValTok{100000}\NormalTok{)}
\end{Highlighting}
\end{Shaded}

Df

SumOfSqs

R2

F

Pr(\textgreater F)

9

257.29

0.37

2.35

0.04023

1

60.20

0.09

4.95

0.04811

1

7.82

0.01

0.64

0.444586

1

13.75

0.02

1.13

0.297517

1

14.72

0.02

1.21

0.280637

1

11.13

0.02

0.92

0.352996

1

73.32

0.11

6.03

0.01695

21

255.22

0.37

36

693.44

1.00

Df

SumOfSqs

R2

F

Pr(\textgreater F)

1

42.57

0.06

3.5

0.070909

1

10.00

0.01

0.82

0.378716

1

19.52

0.03

1.61

0.208418

1

99.98

0.14

8.23

0.00579

1

29.88

0.04

2.46

0.120689

1

140.15

0.20

11.53

0.00173

9

96.14

0.14

0.88

0.563294

21

255.22

0.37

36

693.44

1.00

\newpage

\section{How does species richness relate to network
structure?}\label{how-does-species-richness-relate-to-network-structure}

Following the above results showing the effects of bark roughness and
species richness on network structure, an examination of the variaince
in lichen network similarity explained by species richness and the
network metrics of size and centrality, supported the hypothesis that
network structure is both a function of increasing numbers of species
(i.e., nodes in the network) and other structural differences. In
analyzing two PerMANOVAs with these factors, we found that network size
and centrality both explained vaiation in network similarity regardless
of whether species richness was the first or last factor entered into
the model. Additionally, although variation in network similarity
explained by centrality was low, it was significant even after
partitioning the variance explained by network size. Based on this, we
can conclude that network similarity is both of function of the number
of lichen species in the community determining network size and other
aspects of network strcutre, such as centralization.

\begin{Shaded}
\begin{Highlighting}[]
\NormalTok{cn.l.cen.sr  }\OtherTok{\textless{}{-}}\NormalTok{ vegan}\SpecialCharTok{::}\FunctionTok{adonis2}\NormalTok{(}
\NormalTok{                           cn.d.onc}\SpecialCharTok{\textasciitilde{}}\NormalTok{ L }\SpecialCharTok{+}\NormalTok{ Cen }\SpecialCharTok{+}\NormalTok{ SR,}
                           \AttributeTok{by =} \StringTok{"term"}\NormalTok{, }
                           \AttributeTok{data =}\NormalTok{ onc.dat, }
                           \AttributeTok{mrank =} \ConstantTok{TRUE}\NormalTok{,}
                           \AttributeTok{permutations =} \DecValTok{100000}\NormalTok{)}

\NormalTok{cn.sr.l.cen  }\OtherTok{\textless{}{-}}\NormalTok{ vegan}\SpecialCharTok{::}\FunctionTok{adonis2}\NormalTok{(}
\NormalTok{                           cn.d.onc}\SpecialCharTok{\textasciitilde{}}\NormalTok{ SR }\SpecialCharTok{+}\NormalTok{ L }\SpecialCharTok{+}\NormalTok{ Cen,}
                           \AttributeTok{by =} \StringTok{"term"}\NormalTok{, }
                           \AttributeTok{data =}\NormalTok{ onc.dat, }
                           \AttributeTok{mrank =} \ConstantTok{TRUE}\NormalTok{,}
                           \AttributeTok{permutations =} \DecValTok{100000}\NormalTok{)}


\NormalTok{cn.geno.sr.l.cen  }\OtherTok{\textless{}{-}}\NormalTok{ vegan}\SpecialCharTok{::}\FunctionTok{adonis2}\NormalTok{(}
\NormalTok{                                cn.d.onc}\SpecialCharTok{\textasciitilde{}}\NormalTok{ geno }\SpecialCharTok{+}\NormalTok{ SR }\SpecialCharTok{+}\NormalTok{ L }\SpecialCharTok{+}\NormalTok{ Cen,}
                                \AttributeTok{by =} \StringTok{"term"}\NormalTok{, }
                                \AttributeTok{data =}\NormalTok{ onc.dat, }
                                \AttributeTok{mrank =} \ConstantTok{TRUE}\NormalTok{,}
                                \AttributeTok{permutations =} \DecValTok{100000}\NormalTok{)}

\NormalTok{cn.sr.l.cen.geno  }\OtherTok{\textless{}{-}}\NormalTok{ vegan}\SpecialCharTok{::}\FunctionTok{adonis2}\NormalTok{(}
\NormalTok{                           cn.d.onc}\SpecialCharTok{\textasciitilde{}}\NormalTok{ SR }\SpecialCharTok{+}\NormalTok{ L }\SpecialCharTok{+}\NormalTok{ Cen }\SpecialCharTok{+}\NormalTok{ geno,}
                           \AttributeTok{by =} \StringTok{"term"}\NormalTok{, }
                           \AttributeTok{data =}\NormalTok{ onc.dat, }
                           \AttributeTok{mrank =} \ConstantTok{TRUE}\NormalTok{,}
                           \AttributeTok{permutations =} \DecValTok{100000}\NormalTok{)}
\end{Highlighting}
\end{Shaded}

Df

SumOfSqs

R2

F

Pr(\textgreater F)

1

601.44

0.87

376.37

1e-05

1

36.58

0.05

22.89

3e-05

1

2.69

0.00

1.68

0.195228

33

52.73

0.08

36

693.44

1.00

Df

SumOfSqs

R2

F

Pr(\textgreater F)

1

77.83

0.11

48.71

1e-05

1

536.36

0.77

335.64

1e-05

1

26.52

0.04

16.6

0.00012

33

52.73

0.08

36

693.44

1.00

Df

SumOfSqs

R2

F

Pr(\textgreater F)

9

257.29

0.37

16.95

1e-05

1

46.66

0.07

27.66

2e-05

1

331.43

0.48

196.49

1e-05

1

17.57

0.03

10.42

0.00224

24

40.48

0.06

36

693.44

1.00

Df

SumOfSqs

R2

F

Pr(\textgreater F)

1

77.83

0.11

46.14

1e-05

1

536.36

0.77

317.97

1e-05

1

26.52

0.04

15.72

0.00029

9

12.25

0.02

0.81

0.625554

24

40.48

0.06

36

693.44

1.00

\end{document}
